% Native Chinese reconstruction from SGA4 Springer tome1 scan.
% Source SHA-256 f4f025b21f5fe3c2ef6a94ef525ae510e58b26ecade9712821f8e091df20057b.
% Physical PDF pages 2--13, zero-based indices 1--12. No page overlays.
% SOURCE physical PDF page 2 (page-02.png), visually checked.
\chapter*{第二版序言}

\noindent 1. 正如读者将在本研讨班第一版 (未完成版) 的前言中看到的
(该前言附在本序之后), 研讨班最初的目的是发展\emph{概形的平展上同调}理论.
另一方面, 自口头研讨班举办之年以来, 拓扑和拓扑斯的语言在代数几何中的重要性不断增长:
它既为下降和取商的技术提供方便而直观的框架
(这些技术在几乎所有\emph{概形构造}的问题中都不可或缺),
又用于发展\emph{概形的其他上同调理论},
后者对某些问题比平展上同调更为适用
(例如 fppf 上同调, 或概形的“晶体”上同调).
此外, 现在看来这种语言很可能也将为代数学家和拓扑学家服务.
因此, 越来越需要对这种语言作相当系统的阐述,
并采用尽可能接近其他地方已经通用的记号和术语
(这些记号和术语直接受到拓扑学的启发).
本版除了已经提到的最初目的之外, 也力求给出这样一种阐述,
使其既能作为参考文本, 又能供希望入门拓扑斯语言的数学家使用,
以待将来有更完整的专著问世.
研讨班的标题相对于原来的“概形的平展上同调”有所改变,
% SOURCE physical PDF page 3 (page-03.png), visually checked.
表达的正是这种观点的转变.

出于上述原因, 我们完全重写了原研讨班中讨论拓扑和拓扑斯语言的第 I 至 VI 篇.
我们的指导原则是, 发展一种在各种应用中已经实际使用的语言和记号,
从而不脱离在景之间需要考虑的各种函子的“几何” (或“拓扑”) 内容.
为此, 拓扑斯和拓扑斯态射的概念似乎是不可或缺的主线,
应当置于中心位置, 而景的概念则成为辅助的技术概念.
这尤其促使我们大幅充实了专门讨论这些概念的第 IV 篇,
并依照这一精神完全重写了第 VI 篇 (讨论纤维化景和纤维化拓扑斯).
我们也扩充了讨论范畴和预层的一般概念的第 I 篇,
以便提供研讨班修订形式所需的内部参考.

\noindent 2. 除了相对于研讨班原来形式的这些实质性修改和增补,
还应指出, 本版的另一个区别是加入了关于固有支集上同调和整体对偶定理的第 XVII 和 XVIII 篇.
这两篇由 P. DELIGNE 撰写, 在许多方面与 A. GROTHENDIECK 的口头报告有相当大的区别.
例如, P. DELIGNE 使用 J.L. VERDIER 在拓扑学中的方法,
% SOURCE physical PDF page 4 (page-04.png), visually checked.
直接构造函子 $f^!$, 而不先将态射 $f$ 光滑化;
借助非常有用的“上同调下降理论”, 他还给出固有支集上同调函子 $Rf_!$ 的直接构造,
而不先将分离有限型态射 $f:X\longrightarrow Y$ 紧化;
这后一种理论属于一般拓扑学, 由 B. SAINT-DONAT 撰写为新增的第 $\mathrm{VI}_{\mathrm B}$ 篇,
其在 $Rf_!$ 的直接构造中的应用则由 SAINT-DONAT 在第 XVII 篇的附录中阐述.
另一方面, 我们还要感谢 P. DELIGNE,
他仔细验证了大量相容性, 而他的前任不够谨慎, 对这些相容性曾径直予以承认;
这也使他有机会在第 XVII 篇中相当系统地发展同调代数和导出范畴形式体系的若干补充,
尤其通过系统使用 ind-对象与 pro-对象的技术, 使导出函子的形式体系更加灵活;
因此, 第 XVII 篇的第 1、2、4 节应视为同调代数的标准参考.

\noindent 3. 除了 P. DELIGNE 在阐述和整理方面的这些贡献之外,
研讨班的这次再版还包括了他的一些结果, 主要如下:

\noindent a) 上面已提到的上同调下降理论,
% SOURCE physical PDF page 5 (page-05.png), visually checked.
由 SAINT-DONAT 的第 $\mathrm{VI}_{\mathrm B}$ 篇讨论.
这一理论已经证明是“由局部到整体”的有效工具,
并能 (在可作奇点消解时) 将有奇点概形的某些上同调性质归结到正则概形的情形;
同一原则应当能用于各种解析空间.
(正是如此, P. DELIGNE 最近将它用于发展具有任意奇点的复代数簇的 HODGE 理论.)

\noindent b) 一些“一般拓扑学”的结果, 即关于拓扑斯的结果, 特别是:
局部凝聚拓扑斯存在足够多的点 (VI 9);
拓扑斯上的模的平坦性在逆像函子下保持 (V 8),
这也使他有机会发展一种有用的“局部归纳极限”技术;
以及拓扑斯的 2-纤维积的存在性 (IV 15).

\noindent c) 一个“对称 Künneth 公式” XVII 5.5, 2.1,
借助函子 TS (对称张量积) 的导出函子,
给出对称幂的平展上同调或凝聚上同调;
它受到拓扑学中熟知的公式 (由 DOLD 给出) 的启发.
与可能的猜想相反, 平展上同调中的这个公式\emph{并非}来自“一般抽象论证”,
而是通过逐步归约, 归结到代数闭域上一条固有光滑曲线上的常系数情形;
% SOURCE physical PDF page 6 (page-06.png), visually checked.
在此情形下, 一个特殊化论证和 Lefschetz 原理使问题归结到超越方法的情形,
此时公式由所考虑空间的可三角剖分性得到.
P. DELIGNE 建立对称 Künneth 公式,
是为了将它用于 SGA 5 原有各篇未涵盖的一些情形中的 L 函数公式
(这些内容无疑将收入 SGA 5 的再版),
特别是系数的基环不是域 (如 $\mathbb F_\ell$ 或 $\mathbb Q_\ell$ 这样的域),
或者是素域 $\mathbb F_p$ 的情形 ($p$ 为簇的基域的特征).

\noindent d) 一个“挠子对称幂”的理论 XVII 6.3,
旨在对有限局部自由态射 $f:X\longrightarrow Y$ 定义挠子的迹,
所用条件比通常的条件 ($Y$ 上的群 $G$ 光滑, 或 $f$ 平展) 更为一般.

\noindent e) 一个“挠子积分”的理论 XVIII 1.3,
以及一个“泛系数”公式; 它们作为关于连续系数的对偶理论的组成部分来介绍,
相对于一个光滑且相对维数为 1 的态射 $f:X\longrightarrow S$.
指出 P. DELIGNE 引入的概念启发了平展上同调中 Riemann--Roch 型公式的种种精化:
在由可构造平展层组成的适当的“稳定对象”范畴中写出\emph{典范}同构,
并提示可能向高维推广, 从而开辟了一个有趣的研究领域.
还要指出, 作为技术工具,
P. DELIGNE 对某些光滑相对概形上的某些挠子证明了一个有趣的结构定理;
% SOURCE physical PDF page 7 (page-07.png), visually checked.
这些相对概形 $f:X\longrightarrow S$ 特别包括投射纤维丛 (XVIII 1.2.2),
该定理也提出一些有趣的问题,
尤其与平坦群概形下的挠子的形变理论有关.

\begin{flushright}
Bures, 1969 年 11 月
\end{flushright}

\begin{center}
\textbf{导读}

% Native reconstruction of the LEITFADEN on physical page 7.
% Historical designation VI_B is retained exactly, not silently renumbered.
\begin{tikzpicture}[x=1.35cm,y=1.35cm,font=\small,
  guide/.style={-{Latex[length=2.3mm,width=2.0mm]},double,double distance=1.1pt,line width=.35pt}]
\node (I) at (0,0) {I};
\node (II) at (0,-.35) {II};
\node (III) at (0,-.7) {III};
\node (IV) at (0,-1.05) {IV};
\node (V) at (0,-1.4) {V};
\node (VI) at (0,-1.75) {VI};
\node (VIb) at (1.3,-2.0) {$\mathrm{VI}_{\mathrm B}$};
\draw[decorate,decoration={brace,mirror,amplitude=5pt}] (-.28,.16)--(-.28,-1.91);
\node[align=right,anchor=east] at (-.6,-.88) {拓扑斯理论};
\draw[guide] (VI.east)--(VIb.west);
\node (VII) at (0,-2.5) {VII};
\node (VIII) at (0,-2.85) {VIII};
\node (IX) at (0,-3.2) {IX};
\node (X) at (0,-3.55) {X};
\draw[guide] (VI.south)--(VII.north);
\draw[decorate,decoration={brace,mirror,amplitude=5pt}] (-.28,-2.32)--(-.28,-3.72);
\node[align=right,anchor=east] at (-.6,-2.86) {平展上同调\\一般理论};
\node (XI) at (-3.3,-3.55) {XI};
\draw[guide] (X.west)--(XI.east);
\node (XII) at (0,-4.35) {XII};
\node (XIII) at (0,-4.7) {XIII};
\node (XIV) at (0,-5.05) {XIV};
\draw[guide] (X.south)--(XII.north);
\draw[decorate,decoration={brace,mirror,amplitude=5pt}] (-.28,-4.17)--(-.28,-5.22);
\node[align=right,anchor=east] at (-.6,-4.7) {固有态射\\基变换};
\node (XV) at (-3.3,-4.65) {XV};
\node[align=left,anchor=north west] at (-3.45,-4.84) {光滑态射的\\局部无环性};
\draw[guide] (X.south west)--(XV.north east);
\node (XVI) at (0,-6.0) {XVI};
\draw[guide] (XIV.south)--(XVI.north);
\draw[guide] (-1.95,-5.24)--(XVI.west);
\node (XVII) at (2.1,-6.0) {XVII};
\node[align=left,anchor=west] at (2.4,-6.08) {固有支集上同调\\Künneth};
\draw[guide] (XIV.east)--(XVII.north west);
\node (XIX) at (-3.3,-7.15) {XIX};
\node[align=left,anchor=north west] at (-3.45,-7.35) {特征 0 的\\卓越概形};
\draw[guide] (XVI.south west)--(XIX.north east);
\node (XVIII) at (0,-7.2) {XVIII};
\node[anchor=north] at (0,-7.4) {整体对偶};
\draw[guide] (XVI.south)--(XVIII.north);
\draw[guide] (XVII.south west)--(XVIII.north east);
\end{tikzpicture}
\end{center}

% SOURCE physical PDF page 8 (page-08.png), visually checked.
\chapter*{前言}

\begingroup
\renewcommand{\thefootnote}{\fnsymbol{footnote}}
\noindent 1. 本研讨班的主要目的是发展概形的“WEIL 上同调”的形式体系.
基本上从这里证明的结果出发, 由 WEIL 本人给出的熟知论证,
可以推导出关于有限域上非奇异投射簇的 L 函数的部分 WEIL 猜想
\footnote[1]{写下这些文字时, 尚未证明作用在 $H^i(X,\ZZ_\ell)$ 上的 Frobenius 同态的特征值为代数整数,
更未证明它们的绝对值等于 $q^{i/2}$.}
\footnote[2]{(1968 年 10 月增补).
关于 WEIL 猜想当时进展的综述, 见 KLEIMAN 的报告
\textit{Algebraic cycles and the Weil conjectures},
收入 J. GIRAUD et Al., \textit{Dix exposés sur la cohomologie des schémas}, North Holland Pub. Cie.
[1969 年 8 月增补]. 特征值为代数整数这一事实最近已由 P. DELIGNE 证明
(见 SGA 7 XXI 5).}.
这一推导将在接续本研讨班的一个研讨班中给出 (SGA 5 XIII).
\par\endgroup

本研讨班只研究相对于\emph{平展拓扑}的概形上同调.
从描述上看, 这一拓扑与通常拓扑流形的拓扑十分接近;
我们将看到, 关于通常拓扑空间上同调的大多数经典结果
(各种谱序列、有限性定理、Künneth、对偶、Lefschetz 定理)
都能在新的背景中表述并证明,
只须在必要时限制到与所考虑概形的\emph{剩余特征互素}的挠层.
% SOURCE physical PDF page 9 (page-09.png), visually checked.
通过一个本质上平凡的取投射极限过程 (见 SGA 5 VI),
就能得到 (如 WEIL 所要求的)“特征 0 系数”的上同调理论:
由系数 $\ZZ/\ell^\nu\ZZ$, $\nu\longrightarrow+\infty$,
定义以 $\ell$-进整数环 $\ZZ_\ell$ 为系数的上同调.
当 $\ell$ 与剩余特征互素时,
这种上同调具有经典的 $\ZZ$ 系数上同调中一切通常的良好性质
(因而适于表述 WEIL 猜想).
此外, 当所考虑的概形在复数域 $\CC$ 上为有限型时,
我们将看到 (挠系数) 上同调与相应解析空间的经典上同调本质相同.
一方面, 这将使纯代数方法得到的结果可用于 $\CC$ 上定义的代数簇的通常上同调,
并由此推广各种经典结果;
这些结果通常在非奇异条件下用超越方法证明,
有时还能借此简化某些证明 (尤其在 Lefschetz 理论中).
另一方面, 这还将使我们在适当时,
把超越方法的结果用于特征 $p$ 下可“提升”到特征 0 的代数簇的上同调.

相反, 特征 $p>0$ 所独有的、涉及 $p$-挠系数的本质上新的现象
(例如它们应在局部或整体类域论中起重要作用),
几乎完全超出了
% SOURCE physical PDF page 10 (page-10.png), visually checked.
本笔记采用的平展上同调观点.
它们的研究无疑需要使用更细的拓扑,
很可能是“fppf 拓扑”; 对它的研究此前仅由 SCHATZ 在一个很特殊的情形中开展
[\textit{Cohomology of artinian group schemes over local fields},
Annals of Math. Vol 79, n\textsuperscript{o}3, May 1964, p.411--449]
\begingroup\renewcommand{\thefootnote}{\fnsymbol{footnote}}
\footnote[1]{另见 A. GROTHENDIECK, \textit{le groupe de Brauer III} §11,
收入 \textit{Dix exposés sur la cohomologie des schémas} (North Holland Pub. Cie).}.
在目前看来, 这似乎是平展上同调理论最迫切需要的扩展
\footnote[2]{自写下这些文字以来, 已发现 fppf 拓扑并不满足所期望的一切性质,
而“最迫切的扩展”倒应是发展“晶体上同调”:
它始于 A. GROTHENDIECK 的 \textit{Crystals and the De Rham cohomology of schemes},
受到 MONSKY-WASHNITZER 和 MANIN 工作的启发.}.
还应指出, 需要将 $\pi_i$ 之类同伦不变量的研究推广到概形
(既从平展拓扑的观点, 也从 fppf 拓扑的观点);
可以希望, 将来有待表述的 Hurewicz 型定理,
结合我们对上同调和基本群已有的相当充分的认识,
能使我们沿这条道路取得结果
\footnote[3]{参见原书下一页的\hyperlink{frontprefaces.homotopy-note}{注 (*)}.}.
\par\endgroup

\noindent 2. 除了 WEIL 思想的明显影响之外,
还应指出本研讨班与各种前期研究的联系;
这些研究曾影响我们当中的某些人:

\noindent a) J. TATE 的域的上同调维数理论 [CG],
它本来就受到平展上同调观点的推动,
% SOURCE physical PDF page 11 (page-11.png), visually checked.
而又反过来成为概形的上同调维数理论的一个工具和范例.

\begingroup\renewcommand{\thefootnote}{\fnsymbol{footnote}}
\footnotetext[1]{\hypertarget{frontprefaces.homotopy-note}{}上面提到的研究“平展同伦型”的问题,
在本引言写成以后, 原则上已由 Michael ARTIN 和 Barry MAZUR 两位先生最近的工作解决
(见他们 1966 年关于这一课题的研讨班,
出版为 M. ARTIN et B. MAZUR, \textit{Etale homotopy},
Lecture Notes n\textsuperscript{o}100, Springer (1969)).}

\noindent b) IGUSA 关于“消灭圈”的工作
(它们将在以后的一个研讨班中讨论\footnote[2]{见 SGA 7 III, IX\ldots}),
以及 Igusa 在任意特征下对 \emph{Picard 不等式} $\rho\leq b_2$ 的证明;
一旦有了一个明确的上同调理论,
我们在这里就能相当形式地得到这个不等式
(同时也得到 Néron--Sévéri 的有限性定理)\footnote[3]{见 SGA 6 XIII 5.2, 5.3.}.

\noindent c) OGG 和 CHAFAREVITCH 关于一元函数域以阿贝尔簇为系数的上同调的结果.
这些作者特别得到了代数曲线的\emph{对偶定理}的一个等价形式,
以及一个 Euler--Poincaré 特征公式.
(用平展上同调语言所作的阐述, 见 M. RAYNAUD,
Sém. Bourbaki, 1964 年 12 月\footnote[4]{已重印于原书第 (ix) 页注中所引的文集中.}.)
\par\endgroup

% SOURCE physical PDF page 12 (page-12.png), visually checked.
\noindent 3. 第 I 至 VI 篇从景, 或更确切地说从拓扑斯的观点重新讨论层论,
在许多方面补充了 ARTIN 1962 年 (Harvard) 的研讨班.
其中的阐述以及 ARTIN 中没有的结果, 主要归功于 J. GIRAUD 和 J.L. VERDIER.

第 VII 和 VIII 篇发展关于平展拓扑和平展上同调的最初定义与性质.

第 IX、X 篇给出与挠系数有关的最初的初等定量结果,
尤其涉及代数曲线的上同调、上同调维数以及可构造层的概念.

\begingroup\renewcommand{\thefootnote}{(\arabic{footnote})}
第 XII 至 XVI 篇的结果\footnote[2]{由 M. ARTIN 和 A. GROTHENDIECK 于 1962 年 9 月至 1963 年 3 月间得到.}
围绕两个“基变换定理”展开,
从技术上说, 它们构成本研讨班的中心结果.
这两个定理各自的关键在于基本群的某些性质;
在我们的背景下, 这些性质最自然地用以未必交换的群层为系数的 1-上同调来表述.
因此, 在陈述主要的上同调结果时, 我们也系统地考虑非交换上同调,
这样就能以一种有时更一般、在技术上更方便的形式,
重新得到 SGA 1961 中关于基本群的主要结果
\begingroup\renewcommand{\thefootnote}{\fnsymbol{footnote}}
\footnote[1]{见 Mme M. RAYNAUD 在 SGA 1 XIII 中的报告.}\endgroup.
\par\endgroup

% SOURCE physical PDF page 13 (page-13.png), visually checked.
\begingroup\renewcommand{\thefootnote}{(\arabic{footnote})}
第 XVII 和 XVIII 篇\footnote[3]{由 A. GROTHENDIECK 于 1963 年 2--3 月间得到,
同时得到的还有“局部对偶定理”, 将在 SGA 5 中阐述.}
介绍“紧支集”上同调的理论以及\emph{整体对偶定理},
它们蕴含非奇异代数簇的 \emph{Poincaré 对偶}.

最后, 第 XIX 篇介绍一些围绕“上同调纯性定理”的局部结果
\footnote[4]{由 ARTIN 于 1964 年得到.};
其中大部分使用奇点消解,
所以在有进一步结果之前, 它们只能用于特征零的情形.
\par\endgroup

\noindent 4. 我们将经常使用 EGA I 至 IV, 以及 SGA 1 和 SGA 2 的结果.
